% ===================================================================
%  CADRO N.º 3 — A nenez e a mocidade.
%  Traduction 2026 d'apres texto/io, controlee sur texto/fr, texto/es
%  et texto/pt. Consigne : voir texto/gl/10-cadro-01.tex.
%
%  LES MOIS ET LES JOURS SONT LE PREMIER LIEU OU LE GALICIEN SE
%  SEPARE DE L'ESPAGNOL SANS REJOINDRE LE PORTUGAIS, et ce tableau
%  les nomme tous. Les mois sont latins des trois cotes, mais les
%  jours ne le sont pas : le portugais compte — « segunda-feira,
%  terça-feira » — quand l'espagnol et le galicien nomment.
%  Seulement, le galicien ne nomme pas comme Madrid : « luns »,
%  « martes », « mércores », « xoves », « venres », contre « lunes,
%  martes, miércoles, jueves, viernes ». Quatre formes sur cinq
%  different, et aucune ne ressemble au portugais.
%
%  L'ENTREE DE JUILLET EST DECLAREE DANS gravuri/korekti.json. Le
%  fac-simile ido compose « (julio od agosto) » avec « agosto » en
%  gras et « julio » en romain, alors que les deux sont des mois : la
%  page de lecture met « xullo » en gras comme les autres colonnes.
%  LA SOURCE NE BOUGE PAS ; c'est la lecture qui se corrige, et elle
%  se corrige langue par langue, ce qui fait ici la vingt-sixieme
%  entree de la table.
% ===================================================================

%%K t03-apar-1 apar
\VUsaut{3.0mm}
\VUtitre{15.0pt}{60}{CADRO N.\textsuperscript{o} 3}
\VUsaut{1.6mm}
\VUtitre{11.4pt}{0}{A nenez e a mocidade.}
\VUsaut{3.4mm}

%%K t03-apar-2 apar
(Os cadros 3.º e 4.º están combinados de xeito que presentan, en catro \VUgras{escenas
familiares}, as nocións esenciais sobre o \VUgras{tempo}, a \VUgras{idade}, a
\VUgras{familia}, o \VUgras{vestido} e o \VUgras{estado}.)
\VUsaut{4.0mm}

%%K t03-tit-1 sub
\VUcentre{12.0pt}{0}{\textit{Primeira escena.}}
\VUsaut{3.0mm}

%%K t03-tit-2 sub
\VUpk{10.2pt}{40}{Cerimonia de bautizo nunha aldea.}
\VUsaut{2.4mm}

%%K t03-tit-3 sub
\VUpk{10.2pt}{40}{A primavera.}
\VUsaut{3.0mm}

%%K t03-c1-01-1 p
1. --- Estamos na \VUgras{primavera}, nos meses de \VUgras{marzo}, \VUgras{abril} ou
\VUgras{maio}, nos días da \VUgras{semana}, \VUgras{luns}, \VUgras{martes} ou
\VUgras{mércores}. Son as \VUgras{dez horas} da mañá.

%%K t03-c1-02-1 p
2. --- Fai bo \VUgras{tempo}, o \VUgras{aire} é puro. O sol brila. Algunhas
\VUgras{nubeciñas} \textsuperscript{(2)} soben polo \VUgras{ceo} azul
\textsuperscript{(1)}.

%%K t03-c1-03-1 p
3. --- As \VUgras{árbores} apenas teñen \VUgras{follas}. Os \VUgras{chopos} delgados
\textsuperscript{(3)} e o \VUgras{plátano} alto \textsuperscript{(17)} só teñen ata
agora gomos.

%%K t03-c1-04-1 p
4. --- As \VUgras{andoriñas} \textsuperscript{(4)} volven e voan cun piar ledo
arredor da \VUgras{frecha} \textsuperscript{(7)} do \VUgras{campanario}
\textsuperscript{(5)} que coroa a \VUgras{igrexa} \textsuperscript{(6)} da aldea.

%%K t03-tit-4 sub
\VUsaut{3.4mm}
\VUpk{10.2pt}{40}{A igrexa.}
\VUsaut{3.0mm}

%%K t03-c1-05-1 p
5. --- Moi sinxela é esa igrexa co seu gran \VUgras{pórtico} \textsuperscript{(13)} e
o seu groso \VUgras{reloxo} \textsuperscript{(8)}. A \VUgras{agulla pequena}
\textsuperscript{(9)} está sobre a cifra X; marca as \VUgras{horas.} A
\VUgras{grande} \textsuperscript{(10)} está sobre o XII (mediodía); marca os
\VUgras{minutos.}

%%K t03-c1-06-1 p
6. --- No campanario as \VUgras{campás} \textsuperscript{(11)} tocan ledamente. No
alto do tellado érguese unha pequena \VUgras{cruz} \textsuperscript{(12)} de pedra.

%%K t03-c1-07-1 p
7. --- A igrexa non ten só un gran \VUgras{pórtico} \textsuperscript{(13)}; ten
tamén unha portiña lateral. Por esa portiña o señor \VUgras{párroco}
\textsuperscript{(15)}, vestido coa súa \VUgras{sotana}, sae da \VUgras{sancristía}
\textsuperscript{(14)} e vai á \VUgras{casa reitoral} \textsuperscript{(19)}.

%%K t03-c1-08-1 p
8. --- Preto del, velaquí a \VUgras{leiteira} \textsuperscript{(24)} co seu
\VUgras{burro} \textsuperscript{(23)}. Volve da vila, onde levou bo leite para os
nenos.

%%K t03-tit-5 sub
\VUsaut{4.0mm}
\VUpk{10.2pt}{40}{A horta das froitas.}
\VUsaut{3.4mm}

%%K t03-c1-09-1 p
9. --- O señor Aliboron (o burro) está moi contento con esa carreira da mañá. Aspira
con delicia o aire embalsamado polo recendo das \VUgras{flores}
\textsuperscript{(20)} que cobren as \VUgras{árbores froiteiras} \textsuperscript{(18)}
da \VUgras{horta} veciña. Ben rosados e brancos están eses \VUgras{pexegueiros}
\textsuperscript{(18)} e \VUgras{albaricoqueiros} \textsuperscript{(19)}, e fan
esperar unha boa colleita.

%%K t03-c1-10-1 p
10. --- Mentres tanto, as pequenas \VUgras{abellas} \textsuperscript{(22)} da
\VUgras{colmea} \textsuperscript{(21)} van alá de rapina recoller o seu doce
\VUgras{mel} zoando.

%%K t03-c1-11-1 p
11. --- Pero que pasa? Velaquí que os rapaces da aldea acoden correndo e fan fuxir os
\VUgras{gansos} \textsuperscript{(25)} espantados. É a saída da igrexa do
acompañamento despois do \VUgras{bautizo} do fillo do notario.

%%K t03-c1-12-1 p
12. --- O pequeno Xacobe, no seu \VUgras{berce} \textsuperscript{(26)}, esperta co
tocar ledo das campás, e a súa irmá Madalena, que tamén ela quere ver o
acompañamento do bautizo, pídelle á súa nai que veña peitearlle o cabelo e vestila na
soleira da casa.

%%K t03-c1-13-1 p
13. --- A nai consente. Pon o seu \VUgras{pano} \textsuperscript{(28)}, o seu
\VUgras{corpiño} \textsuperscript{(29)} e o seu \VUgras{mandil}
\textsuperscript{(30)}, e vén sentar nunha cadeiriña diante da porta. Madalena só ten
o seu \VUgras{refaixo} \textsuperscript{(31)} e o seu \VUgras{corselete}
\textsuperscript{(32)}.

%%K t03-c1-14-1 p
14. --- Que feliz está! Esquece as súas flores preferidas, os seus \VUgras{caraveis}
\textsuperscript{(34)}, sobre os que pousan as \VUgras{bolboretas}
\textsuperscript{(35)}, as súas \VUgras{tulipas} \textsuperscript{(36)}, os seus
\VUgras{lirios do val} \textsuperscript{(37)}, en \VUgras{testos}
\textsuperscript{(38)}, sobre o \VUgras{muro} \textsuperscript{(33)}, e as súas
\VUgras{rosas} \textsuperscript{(39)}, que fan unha sebe tan fermosa. Ela mira a rica
roupa das damas do acompañamento.

%%K t03-c1-15-1 p
15. --- Os rapaces da aldea non prestan moita atención á roupa. Precipítanse e
empúrranse uns aos outros para acadaren \VUgras{larpeiradas} \textsuperscript{(55)} ou
\VUgras{moedas} \textsuperscript{(52)}. Un deles chora \textsuperscript{(40)}.

%%K t03-c1-16-1 p
16. --- O outro pisou a pata do \VUgras{can} \textsuperscript{(42)}, que foxe a
ouvear e fai fuxir a \VUgras{galiña} \textsuperscript{(43)} e mais os seus
\VUgras{pitiños} \textsuperscript{(44)}.

%%K t03-tit-6 sub
\VUsaut{5.0mm}
\VUpk{10.2pt}{40}{O desfile do bautizo.}
\VUsaut{4.0mm}

%%K t03-c1-17-1 p
17. --- Diante do acompañamento camiña a \VUgras{ama de cría}
\textsuperscript{(45)}. Ten unha bela \VUgras{cofia} \textsuperscript{(46)} con
longas cintas. Leva o \VUgras{recentemente nado} \textsuperscript{(47)}, todo vestido
de \VUgras{encaixes} \textsuperscript{(48)}.

%%K t03-c1-18-1 p
18. --- Detrás da ama veñen o \VUgras{padriño} \textsuperscript{(49)} e a
\VUgras{madriña} \textsuperscript{(53)}. Reparten as larpeiradas e as moedas. O
padriño ten \VUgras{luvas} \textsuperscript{(51)} e \VUgras{chapeu de copa}
\textsuperscript{(50)}. A madriña ten na man un fermoso \VUgras{ramo}
\textsuperscript{(54)}.

%%K t03-c1-19-1 p
19. --- Detrás deles vemos o \VUgras{tío} \textsuperscript{(56)} e a \VUgras{tía}
\textsuperscript{(58)} do neno. A tía ten un elegante \VUgras{sombreiro}
\textsuperscript{(59)}, o tío unha \VUgras{levita} negra \textsuperscript{(57)} e un
chaleco branco. Velaquí despois o \VUgras{curmán} \textsuperscript{(60)} e a
\VUgras{curmá} \textsuperscript{(61)}. Esta última leva pola man a \VUgras{irmá}
\textsuperscript{(62)} do recentemente nado, a pequena Berta.

%%K t03-c1-20-1 p
20. --- O outro \VUgras{irmán} \textsuperscript{(63)} de Berta camiña detrás. Dálle a
man a unha \VUgras{parenta} \textsuperscript{(64)}. Os \VUgras{amigos}
\textsuperscript{(65)} da familia veñen despois.

%%K t03-tit-7 sub
\VUsaut{4.6mm}
\VUpk{10.2pt}{40}{Os que miran.}
\VUsaut{3.4mm}

%%K t03-c1-21-1 p
21. --- Preto do acompañamento, velaquí un \VUgras{pedinte} \textsuperscript{(66)}
apoiado na súa \VUgras{muleta} \textsuperscript{(67)}. Pide esmola (mendiga).

%%K t03-c1-22-1 p
22. --- Do outro lado hai un rapaciño moi \VUgras{ben educado}
\textsuperscript{(68)}. Saúda e dille «\VUgras{bos días}» ás persoas que coñece.

%%K t03-c1-23-1 p
23. --- Os aldeáns saíron das súas casas para veren o desfile do bautizo. Vemos un
\VUgras{labrego} \textsuperscript{(69)} co seu \VUgras{gorro de algodón} e á beira
unha \VUgras{labrega} \textsuperscript{(70)}. Despois unha \VUgras{vella}
\textsuperscript{(71)} apoiada no seu \VUgras{bastón} \textsuperscript{(72)}.
Finalmente o \VUgras{muiñeiro} \textsuperscript{(73)} co seu gran \VUgras{sombreiro
de feltro} \textsuperscript{(74)} e unha labrega nova calzada con \VUgras{zocas}
\textsuperscript{(75)}, que ensina a camiñar ao seu neniño.

%%K t03-c1-24-1 p
24. --- Esa escena é moi divertida.

%%K t03-tit-8 sub
\VUsaut{4.0mm}
\VUcentre{12.0pt}{0}{\textit{Segunda escena.}}
\VUsaut{3.0mm}

%%K t03-tit-9 sub
\VUpk{10.2pt}{40}{Festa pública.}
\VUsaut{2.4mm}

%%K t03-tit-10 sub
\VUpk{10.2pt}{40}{Xogos náuticos e regatas.}
\VUsaut{3.0mm}

%%K t03-c2-01-1 p
1. --- Chegou o \VUgras{verán}. O sol ardente do mes de \VUgras{xuño} (xullo ou
\VUgras{agosto}) anima a mocidade a bañarse e a andar de barca.

%%K t03-c2-02-1 p
2. --- Na ribeira dereita do río, os remadores desvístense para participaren nos xogos
náuticos e nas regatas.

%%K t03-c2-03-1 p
3. --- Un deles quita os seus \VUgras{zapatos} \textsuperscript{(1)}; desátaos
(desabotóaos). Os seus dous compañeiros xa están listos. Desde agora só teñen a súa
\VUgras{camiseta de punto} \textsuperscript{(2)} e o seu \VUgras{calzón de baño}
\textsuperscript{(3)}. Conversan mentres esperan polos amigos. Falan da feira e da
festa que se celebra na outra ribeira.

%%K t03-c2-04-1 p
4. --- No chan, preto deles, xace a \VUgras{roupa} dos seus compañeiros: un
\VUgras{par} de \VUgras{botíns} \textsuperscript{(4)}, \VUgras{zapatos}
\textsuperscript{(5)}, \VUgras{calcetíns} \textsuperscript{(6)}, sombreiros de
\VUgras{feltro} \textsuperscript{(7)} e de \VUgras{palla} \textsuperscript{(8)} e unha
\VUgras{gravata} \textsuperscript{(9)}.

%%K t03-c2-05-1 p
5. --- Un dos mozos dobrou con coidado o seu \VUgras{traxe} \textsuperscript{(10)}.
Pono sobre unha toalla, na que o seu veciño ha meter axiña os dous \VUgras{traxes.}

%%K t03-c2-06-1 p
6. --- O sexto rapaz tarda. Aínda ten a súa \VUgras{pucha} \textsuperscript{(11)} na
cabeza. Non quitou nin a súa \VUgras{camisa} \textsuperscript{(12)}, nin o seu
\VUgras{colo} \textsuperscript{(13)}, nin os seus \VUgras{puños}
\textsuperscript{(14)}. Ten na man o seu \VUgras{chaleco} \textsuperscript{(17)} e
leva aínda os seus \VUgras{pantalóns} \textsuperscript{(16)} e os seus
\VUgras{tirantes} \textsuperscript{(15)}. De certo non estará listo para a carreira
próxima.

%%K t03-c2-07-1 p
7. --- Preto dos mozos, un carreiro nas beiras do cal hai moitos \VUgras{cardos}
\textsuperscript{(18)} leva á beira do río, onde o pai Xacobe prepara entre as
\VUgras{canas} \textsuperscript{(19)} os \VUgras{fogos de artificio}
\textsuperscript{(20)} que se dispararán á noite.

%%K t03-tit-11 sub
\VUsaut{4.4mm}
\VUpk{10.2pt}{40}{A carreira.}
\VUsaut{3.4mm}

%%K t03-c2-08-1 p
8. --- No río algunhas damas, resgardándose cos seus parasoles, pasean nunha
\VUgras{barca} \textsuperscript{(21)}. Seguen a carreira, que é moi interesante. En
cada \VUgras{iola} \textsuperscript{(22)}, catro \VUgras{remeiros}
\textsuperscript{(24)} reman con todas as súas forzas, mentres o \VUgras{timoneiro}
\textsuperscript{(26)} ten o \VUgras{temón} \textsuperscript{(25)} e dirixe o fráxil
bote. Os \VUgras{remos} \textsuperscript{(23)} baten a auga a compás e as lixeiras
embarcacións navegan cunha rapidez vertixinosa.

%%K t03-c2-09-1 p
9. --- Algúns mozos loitan entre si, preto do piar da ponte, polo premio do
\VUgras{mastro de gurupés} \textsuperscript{(28)}. Un deles camiña con prudencia sobre
o mastro flexible e resvaladío para acadar a pequena \VUgras{bandeira}
\textsuperscript{(29)} fixada na punta. O seu desafortunado \VUgras{competidor}
\textsuperscript{(30)} caeu na auga. Nada para chegar a terra.

%%K t03-c2-10-1 p
10. --- Mozos de máis idade \VUgras{xustan en barca} \textsuperscript{(32)} preto do
\VUgras{peirao} \textsuperscript{(33)}. A todos eses xogos asiste un \VUgras{público}
\textsuperscript{(31)} numeroso apoiado no \VUgras{parapeto} da \VUgras{ponte}
\textsuperscript{(27)} e amoreado na \VUgras{ribeira.} Algúns espectadores, porén,
volven a cabeza para seguiren o xogo do \VUgras{mastro do premio}
\textsuperscript{(45)} ou para admiraren o \VUgras{acompañamento do alcalde} que vén
para o \VUgras{reparto} dos premios.

%%K t03-c2-11-1 p
11. --- Primeiro, velaquí algúns \VUgras{cativos} \textsuperscript{(35)} que van
diante dos \VUgras{músicos} \textsuperscript{(36)}; despois vén o \VUgras{alcalde}
\textsuperscript{(37)}, os tenentes de alcalde e o \VUgras{concello} da vila.
Finalmente camiñan os \VUgras{bombeiros} \textsuperscript{(38)}.

%%K t03-tit-12 sub
\VUsaut{5.0mm}
\VUpk{10.2pt}{40}{A feira.}
\VUsaut{3.6mm}

%%K t03-c2-12-1 p
12. --- Hai moitísima xente na \VUgras{praza da feira} \textsuperscript{(39)}.
Vén a xente ver as diversas \VUgras{barracas}: os \VUgras{cabalos de madeira}
\textsuperscript{(40)}, o \VUgras{circo} \textsuperscript{(41)}, no que xa comezou a
función; o \VUgras{baile público} \textsuperscript{(42)}, onde os mozos e as mozas da
vila gozan do pracer do \VUgras{baile.}

%%K t03-c2-13-1 p
13. --- Ao fondo vemos a \VUgras{praza de touros} \textsuperscript{(44)}, na que o
valente \VUgras{matador} español loita contra o furioso \VUgras{touro}
\textsuperscript{(45)}, despois o \VUgras{tobogán} \textsuperscript{(49)}, a
\VUgras{barraca de tiro} \textsuperscript{(46)} onde se exercita un a usar as
\VUgras{armas de fogo} e a tirar ao \VUgras{branco.}

%%K t03-c2-14-1 p
14. --- Preto de alí, velaquí o \VUgras{teatro de feira} \textsuperscript{(47)}, no
que se representan a \VUgras{comedia} e o \VUgras{drama.} Moi preto está a barraca do
\VUgras{xigante} \textsuperscript{(48)} e do \VUgras{anano.} A \VUgras{estatura} do
primeiro é de dous metros e quince centímetros; o segundo é apenas tan alto coma un
neno.

%%K t03-c2-15-1 p
15. --- Finalmente, velaquí a gran \VUgras{casa de feras} \textsuperscript{(50)} coas
súas \VUgras{feras}, o seu \VUgras{tigre} \textsuperscript{(51)} de poderosas
\VUgras{unllas}, o seu \VUgras{león} \textsuperscript{(52)} de mesta \VUgras{melena},
as súas dúas \VUgras{xirafas} \textsuperscript{(53)} e o seu \VUgras{elefante}
\textsuperscript{(54)}, que usa con tanta destreza a súa \VUgras{trompa.}

%%K t03-c2-16-1 p
16. --- O \VUgras{domador} \textsuperscript{(55)} que amestra esas feras está á porta
da barraca.
